\section{A positive finite-field word for the binary endpoint tensor}
\label{app:positive-word}

The address field $\mathbb F_3$ above is not the phase field. For the
constant word use
\[
 E=\mathbb F_2[t]/(t^4+t+1),\qquad \eta=t^5=t^2+t.
\]
The quartic has no root in $\mathbb F_2$ and remainder $1$ modulo the
only irreducible quadratic $t^2+t+1$, so it is irreducible. Thus
$E=\mathbb F_{16}$, it contains $\mathbb F_4=\mathbb F_2(\eta)$, and
$\eta^2+\eta+1=0$. The four-state phase letter is
\[
 A=(\eta I_3+\eta^2J_3)\oplus[1]=C\oplus[1].
\]
It acts separately on the two $\mathbb F_4$ coefficient planes in
$E=\mathbb F_4\oplus t\mathbb F_4$. Its inverse is $A^2$.
The inverse-phase wrapper on these planes is
\begin{equation}
 \widehat\sigma(a+tb)=a^2+t b^2\quad(a,b\in\mathbb F_4),       \label{eq:positive-17}
\end{equation}
which applies the $\mathbb F_4$ Frobenius to each plane and fixes the
codec basis element $t$. It is $\mathbb F_2$-linear and is not ordinary
$\mathbb F_{16}$ squaring.
In the polynomial basis $(1,t,t^2,t^3)$ it is the explicit map
\[
 (c_0,c_1,c_2,c_3)\longmapsto
 (c_0+c_2+c_3,\ c_1+c_3,\ c_2,\ c_3).
\]
Indeed $a_0=c_0$, $a_1=c_2+c_3$, $b_0=c_1+c_2+c_3$ and
$b_1=c_3$ in $a=a_0+a_1\eta$, $b=b_0+b_1\eta$.
Three in-place XORs implement it, and repeating them gives the identity.
For $\lambda\in\F_4$,
$\widehat\sigma(\lambda z)=\lambda^2\widehat\sigma(z)$.
Consequently coefficientwise conjugation gives
\[
 \widehat\sigma C\widehat\sigma
   =\eta^2I_3+\eta J_3=C^2=C^{-1},
 \qquad
 \widehat\sigma(C^{\otimes f})\widehat\sigma
   =(C^{-1})^{\otimes f}.
\]
This holds on arbitrary $E$-valued inputs. The two passes cancel on
idle or ineligible coordinates; either pass alone can change their
coefficient values. Both passes act within the existing four-bit lanes,
in $O(V)$ time around one child. Ordinary $E$ squaring used on both
sides would fail even on an idle value $t$, sending it to
$t^4=t+1$ instead of $t$. The semilinearity assertion is only for
$\F_4$ scalars, not arbitrary $E$ scalars.

Here is the explicit word construction. Put $a=\eta^2$ and $b=\eta$;
use column vectors and rightmost-first composition. Then
\[
 A^3=I,\quad A^{-1}=A^2,\quad
 u=(1,a,a,0)^T,\quad v=(1,b,b,0)^T,\quad
 P=AE_{11}A^{-1}=uv^T,
\]
with $v^Tu=1$, hence $P^2=P$. Let
\[
 D=\operatorname{diag}(t,1,t+1,1),\quad Q=DPD^{-1}=x_0y_0^T,
\]
where $x_0=(t,a,a(t+1),0)^T$ and
$y_0=(t^{-1},b,b/(t+1),0)^T$. Direct products give
\begin{equation}
 v^Tx_0=0,\qquad y_0^Tu=\kappa
 =\frac{t^2+t+1}{t(t+1)}\ne0,
 \qquad PQ=0,\quad QP=\kappa x_0v^T.                 \label{eq:positive-18}
\end{equation}
For idempotents $P,Q$ with $PQ=0$,
\begin{equation}
 [I+cP,I+dQ]=I+\frac{cd}{1+c}QP.                    \label{eq:positive-19}
\end{equation}
This follows by using $(I+cP)^{-1}=I+c(1+c)^{-1}P$ and
$QPQ=0$. With $c=d=\eta^2$, $cd/(1+c)=1$. Thus
\begin{equation}
\begin{aligned}
 G&=A\operatorname{diag}(\eta,1,1,1)A^{-1}=I+\eta^2P,\\
 H&=DGD^{-1}=I+\eta^2Q,\\
 T_0&=GHG^{-1}H^{-1}=I+\kappa x_0v^T.
\end{aligned}\label{eq:positive-20}
\end{equation}
Each of $G,H,G^{-1},H^{-1}$ uses three forward $A$ letters after
replacing inverse $A$ by $A^2$, so $T_0$ uses 12.

The remaining conjugations isolate an elementary matrix. Set
\[
 B=\operatorname{diag}(1,\eta,\eta,1),\quad
 x=(t,1,t+1,0)^T,\quad w=(1,1,1,0)^T.
\]
Then $T_1=BT_0B^{-1}=I+\kappa xw^T$ (12 $A$ letters). Let $P_{12}$
swap coordinates 1 and 2. Since $w$ is fixed and $w^TP_{12}x=0$,
\begin{equation}
 T_2=T_1P_{12}T_1P_{12}^{-1}
   =I+\kappa(t+1)(e_1+e_2)w^T                       \label{eq:positive-21}
\end{equation}
uses 24 $A$ letters. With $D_3=\operatorname{diag}(1,1,\eta,1)$,
\begin{equation}
 T_3=T_2D_3T_2D_3^{-1}
   =I+\kappa(t+1)\eta(e_1+e_2)e_3^T                \label{eq:positive-22}
\end{equation}
uses 48; the two cross pairings vanish and the row sum is
$(1+\eta^{-1})e_3^T=\eta e_3^T$. With
$D_1=\operatorname{diag}(\eta,1,1,1)$,
\begin{equation}
 T_4=T_3D_1T_3D_1^{-1}
   =I+\kappa(t+1)E_{13}                             \label{eq:positive-23}
\end{equation}
uses 96; the cross pairings vanish and the column factor is
$1+\eta=\eta^2$. Put $\mu=\kappa(t+1)\ne0$; then
\begin{equation}
 T_5=\operatorname{diag}(\mu^{-1},1,1,1)
       T_4\operatorname{diag}(\mu,1,1,1)=I+E_{13}.  \label{eq:positive-24}
\end{equation}
All conjugating factors are invertible monomial matrices. Coordinate
permutations produce every $U_{ij}=I+E_{ij}$ with the same 96 forward
$A$ letters. This is an explicit formula, not an uncomputed word search.

Finally write $H_2=US$, with $U=I+E_{12}$ and $S$ the two-coordinate
swap. In ordered coordinates $(00,01,10,11)$, put
\[
 P_{\rm first}=(1\ 3)(2\ 4),\qquad
 P_{\rm second}=(1\ 2)(3\ 4).
\]
Then, using the four explicit elementary-matrix words,
\begin{equation}
 H_2\otimes H_2
 =U_{13}U_{24}P_{\rm first}U_{12}U_{34}P_{\rm second}.       \label{eq:positive-25}
\end{equation}
There are four $U_{ij}$ factors, hence $4\cdot96=384$ forward $A$
letters; the remaining factors expand to 392 invertible monomial
letters, for 776 total. The completed matrix has entries in $\mathbb F_2$.
As an $\mathbb F_2$-linear map on the four binary coordinates of an
$E$-symbol, it applies $H_2\otimes H_2$ independently to each coordinate
plane. Thus a binary $0/1$ embedding returns to $0/1$, and any zero lane
used for padding returns to zero at the completed word. An arbitrary
$E$-valued input need not return to the scalar subfield. Intermediate
$E$ lanes must not be cropped.
When four arbitrary payload bits are packed into one $E$ symbol,
all four coefficient planes remain data; only extra zero bits from
the final partial group are removed. No projection $E\to\F_2$
is applied to such a packed symbol.

Pairing $2h$ identical native endpoint axes and tensoring this fixed word
reduces their algebraic action to $A^{\otimes h}$ letters and fixed
monomial letters. The permutation part of each monomial is in
$\mathrm{AGL}(2,\mathbb F_2)=S_4$: its NOT/CNOT address shears and the
product of its $h$ diagonal $E$ scalars are charged by the separately
supplied native mask/routing and record-scan interfaces. The active-sector
role construction reduces the $A^{\otimes h}$ letters. The native
large-even-radix endpoint appendix supplies their tape price. No binary
$2^k$-axis H2 bound is asserted here; that requires a distinct binary
mask-routing argument and is not needed for the manuscript's native
endpoint interface.
For precision, write a monomial as $M_{\pi(d),d}=m_d$, with
$m_d\in E^*$. Apply the coefficient $\prod_i m_{d_i}$ at its current
\emph{source} address, then route by $\pi^{\otimes h}$. Equivalently
an output-address scan would use $\prod_i m_{\pi^{-1}(e_i)}$.
Using the source-indexed list unchanged after routing is generally
incorrect. Tensoring the fixed word is letterwise:
$(L_1\cdots L_{776})^{\otimes h}
 =L_1^{\otimes h}\cdots L_{776}^{\otimes h}$.
Thus there are 384 tensor phase calls, not 384 separate calls per axis.

